\clearpage
\begingroup
\fancyhead[L]{\small THEORY OF THE MOTION OF THE MOON.}
\renewcommand{\headrulewidth}{0pt}
\SourceLeaf{623}
\section*{SECOND SECTION.}
\begin{center}
\itshape Integration of the fundamental equations, with the perturbing
forces of the Sun set aside.
\end{center}

\par\noindent\hspace*{\parindent}\textls[120]{Theorem from integral calculus:}

\par\noindent\hspace*{\parindent}The integral of the second-order differential equation
\[
0=\frac{\mathrm{dd}w}{\mathrm{d}v^2}+\alpha\alpha w+Z,
\]
\noindent where \(Z\) denotes a function of \(v\) and \(\alpha\alpha\) a constant
positive quantity, is
\[
\begin{aligned}
\alpha w&=\cos\alpha v\int\sin\alpha v\mathpunct{.}\,Z\mathrm{d}v
 -\sin\alpha v\int\cos\alpha v\mathpunct{.}\,Z\mathrm{d}v\\
 &\qquad+\mathit{Const.}\sin(\alpha v+\mathit{Const.}).
\end{aligned}
\]

\par\noindent\hspace*{\parindent}Thus, if \(Z=a+b\sin\beta v+c\sin\gamma v+\mathrm{etc.}\), we have
\[
\begin{aligned}
w&=-\frac{a}{\alpha\alpha}+\frac{b}{\beta\beta-\alpha\alpha}\sin\beta v
 +\frac{c}{\gamma\gamma-\alpha\alpha}\sin\gamma v\\
 &\qquad+\mathrm{etc.}+\mathit{Const.}\sin(\alpha v+C.).
\end{aligned}
\]

\par\noindent Or, if \(Z=a+b\cos\beta v+c\cos\gamma v+\mathrm{etc.}\), we have
\[
\begin{aligned}
w&=-\frac{a}{\alpha\alpha}+\frac{b}{\beta\beta-\alpha\alpha}\cos\beta v
 +\frac{c}{\gamma\gamma-\alpha\alpha}\cos\gamma v\\
 &\qquad+\mathrm{etc.}+\mathit{Const.}\cos(\alpha v+C.).
\end{aligned}
\]

\par\noindent\hspace*{\parindent}If, on the other hand, \(Z\) contained a term \(a\cos\alpha v\), the corresponding term to be put in
\(w\) would be \(-\tfrac12\frac{av\sin\alpha v}{\alpha}
-\tfrac14\frac{a\cos\alpha v}{\alpha\alpha}\), or, for a term
\(a\sin\alpha v\) in \(Z\), the term appearing in \(w\) would be
\(\tfrac12\frac{av\cos\alpha v}{\alpha}
-\tfrac14\frac{a\sin\alpha v}{\alpha\alpha}\mathpunct{\raisebox{.4ex}{.}}\)
Thus, if it is known a priori that \(w\) can contain no circular arcs,
\(Z\) can contain neither a sine nor a cosine of \(\alpha v\).

\par\noindent\hspace*{\parindent}We observe that each of the fundamental equations above consists of two
parts, one of which contains the quantities \(\pi,\omega,\psi\) as a factor,
while the other is independent of them. If the motion of the Moon
were unperturbed by the Sun, the first parts would drop out, and the
inte\SourceLeaf{624}gration\marginpar{\scriptsize VII.\quad78} would then present no difficulty. We also see from the theorem above
that the integration could be carried out even with those first terms retained,
if we could express them as functions of \(v\). But this cannot be done
without already knowing the motion of the Moon.
Fortunately, the effects of the Sun are quite small.
The greatest periodic equation for the longitude of the Moon
amounts to only \(1\tfrac13\) degrees, and the apogee of the Moon, which
would remain fixed without the action of the Sun, acquires from it
only a slow prograde motion, whose ratio to the mean motion of the
Moon is 1 to 118. The retrograde motion of the node is still smaller.
Without these circumstances, calculating the motion
of the Moon would exceed the powers of analysis. These circumstances make it possible
to approach the true motion of the Moon in successive steps. We
first substitute the purely elliptic motion of the Moon in the parts that depend on \(\pi,\omega,\psi\),
and derive a first approximation to the motion from this.
We then use this motion in those same parts in place of the elliptic motion
to obtain a more accurate determination, repeating these successive
improvements as often as the desired degree of accuracy
requires.

\par\noindent\hspace*{\parindent}After these preliminary remarks, we first take up the equations above
with the quantities \(\pi,\omega,\psi\) set aside, in order to
obtain the elliptic motion. In doing so we may also neglect quantities
that depend on higher powers of the eccentricity and
inclination.

\par\noindent\hspace*{\parindent}The third equation here becomes
\[
\frac{\mathrm{dd}\theta}{\mathrm{d}v^2}+\theta=0,
\]
\noindent whose integral is (exactly)
\[
\theta=g\sin\lambda,\quad\text{so that}\quad\lambda=v+\mathit{Const.}
\]
\noindent This shows that the orbit of the Moon, viewed from the center of the Earth,
[is] a great circle whose inclination to the ecliptic has
tangent \(g\). The longitude of the node \(=-\mathit{Const.}\), or, equivalently, \(\lambda\)
is the argument of latitude\HistoricalNoteEight{But reckoned along the ecliptic,
not along the orbit of the Moon.]}.

\SourceLeaf{625}
\par\noindent\hspace*{\parindent}The second equation becomes
\[
0=\frac{\mathrm{dd}u}{\mathrm{d}v\mathrm{d}u}
 +\frac{2\mathrm{d}p}{p\mathrm{d}v},
\]
\noindent from which we get
\[
\log\frac{pp\mathrm{d}u}{\mathrm{d}v}=\log\mathit{Const.},
\]
\noindent hence
\[
\frac{pp\mathrm{d}u}{\mathrm{d}v}=\mathit{Const.}=A.
\]
\noindent The first equation accordingly becomes
\[
0=\frac{\mathrm{dd}p}{\mathrm{d}v^2}+p-AAk(1+\theta\theta)^{-\frac32},
\]
\noindent if for brevity we set \(\frac{(m+m')\mathrm{d}t^2}{\mathrm{d}u^2}=k\).
Hence, substituting the value of \(\theta\) given above and neglecting the higher
powers of \(g\), we obtain
\[
0=\frac{\mathrm{dd}p}{\mathrm{d}v^2}+p
-AAk\left(1-\tfrac34gg+\tfrac34gg\cos2\lambda\right).
\]
\noindent From this we obtain \(p=AAk(1-\tfrac34gg)(1-e\cos M-\tfrac14gg\cos2\lambda)\).
Thus, for the constant part of \(p\), we set \(AAk(1-\tfrac34gg)=h\) and obtain
\[
\begin{aligned}
p&=h(1-e\cos M-\tfrac14gg\cos2\lambda)\\[5pt]
\log p&=\log h-\tfrac14ee-e\cos M-\tfrac14ee\cos2M-\tfrac14gg\cos2\lambda\\[5pt]
\log pp\frac{\mathrm{d}u}{\mathrm{d}v}
 &=\log A=\tfrac12\log h-\tfrac12\log k+\tfrac38gg\\[5pt]
\log\frac{\mathrm{d}u}{\mathrm{d}v}
 &=-\tfrac32\log h-\tfrac12\log k+\tfrac12ee+\tfrac38gg\\
 &\qquad+2e\cos M+\tfrac12ee\cos2M+\tfrac12gg\cos2\lambda\\[5pt]
\frac{\mathrm{d}u}{\mathrm{d}v}
 &=h^{-\frac32}k^{-\frac12}(1+\tfrac32ee+\tfrac38gg)\\
 &\qquad(1+2e\cos M+\tfrac32ee\cos2M+\tfrac12gg\cos2\lambda),
\end{aligned}
\]
\noindent consequently
\[
1+ee+\tfrac14gg=hk^{\frac13}
\]
\noindent and
\[
u=v+2e\sin M+\tfrac34ee\sin2M+\tfrac14gg\sin2\lambda,
\]
\noindent where we have neglected all quantities containing higher powers of
\(e\) and \(g\). Incidentally, here \(e\) is evidently a constant to be determined from
observations (the eccentricity), and \(M\) is the
true anomaly\HistoricalNoteEight{But reckoned along the ecliptic,
not along the orbit of the Moon.]} \(=v-\mathit{Const.}\), so that
\(\mathit{Const.}\) [is] the invariable longitude of the apogee.

\par\medskip\noindent\hfill\rule{.19\textwidth}{.3pt}\hfill\mbox{}\par
\noindent\hfill 78\textsuperscript{*}\par
\clearpage
\endgroup
